\section{Introduction}
Symmetry principles provide exact analytical constraints on quantum and classical mechanical Hamiltonians without requiring the explicit solutions to complex differential equations~\cite{Ramadevi2019, LandauMechanics}. In polyatomic systems, spatial symmetries governed by discrete point groups determine the degeneracy of vibrational eigenstates, identify forbidden quantum transitions, and govern selection rules in spectroscopy~\cite{Cotton2008}.

In this review, we contrast two foundational archetypes of molecular symmetry: the dihedral system $D_{nh}$ (featuring an $n$-fold principal axis and $n$ transverse two-fold axes), and the exceptional non-crystallographic icosahedral system $Y_h$ ($I_h$). While $D_{6h}$ (benzene, $\text{C}_6\text{H}_6$) represents the pinnacle of planar $\pi$-conjugated chemistry, $Y_h$ (buckminsterfullerene, $\text{C}_{60}$) is the highest-order point group physically realized in molecular physics, possessing 120 symmetry operations.
